\chapter{Results and Analysis}
\label{chap:results}

\begin{infobox}
\small\textbf{Validity status.}
All retrieval, baseline and benchmark figures in this chapter are computed against
the \emph{frozen, independently reviewed} benchmark of 98 records described in
Section~\ref{sec:frozen-benchmark}. Generation results (Section~\ref{sec:generation}) are reported on
\textbf{both partitions}: the held-out run was made once, after every component
it evaluates had been frozen, and nothing was revised in response to it.
Escalation calibration (Section~\ref{sec:escalation}) remains
\textbf{development-only}, because the one component it would fit --- the
retrieval-uncertainty threshold --- is deliberately left uncalibrated. Every
figure is reproduced from a stored experiment manifest, metrics file or validation
record; where a quantity exists only in a written report, the text says so.
\end{infobox}

\section{Overview}

This chapter reports what the system does and where it fails.
It proceeds from the data outwards: the corpus and the benchmark
(Section~\ref{sec:frozen-benchmark}), retrieval
(Section~\ref{sec:retrieval-results}), the historical three-system baseline
re-scored under a corrected evaluator (Section~\ref{sec:baselines}),
evidence-constrained generation and its verification
(Section~\ref{sec:generation}), and the calibration of the escalation policy
(Section~\ref{sec:escalation}). Section~\ref{sec:invisible} isolates two
properties of the benchmark that no internal metric of the pipeline can detect,
and Section~\ref{sec:threats} states the threats to validity.

Table~\ref{tab:results-summary} collects the chapter's principal results with a pointer to where each is established.

\begin{table}[htbp]
  \centering
  \caption{Principal results of this chapter, with the section or table in which
  each is reported. Held-out figures come from a single run made after every
  component they evaluate had been frozen.}
  \label{tab:results-summary}
  \small
  \begin{tabularx}{\textwidth}{@{}Xrl@{}}
    \toprule
    \textbf{Quantity} & \textbf{Value} & \textbf{Reported in} \\
    \midrule
    \multicolumn{3}{@{}l}{\textit{Corpus and benchmark}} \\
    Legal parents / retrieval children & 109 / 1{,}535 & \S\ref{sec:frozen-benchmark} \\
    Frozen benchmark after independent review & 98 records & Table~\ref{tab:freeze} \\
    \quad development / held-out partition & 42 / 56 & Table~\ref{tab:freeze} \\
    \quad legal parents shared between partitions & 0 & Table~\ref{tab:freeze} \\
    Reviewer agreement, Fleiss' $\kappa$, three categories & 0.250 & Table~\ref{tab:agreement} \\
    \addlinespace
    \multicolumn{3}{@{}l}{\textit{Retrieval, held-out ($n = 56$)}} \\
    Hybrid with cross-encoder, recall @1 / @5 / @20 & 71.4 / 89.3 / 96.4\% & Table~\ref{tab:retrieval-test} \\
    Depth-20 recall gained by reranking & none & \S\ref{sec:retrieval-results} \\
    Mean evidence recall at depth 5 & 81.1\% & \S\ref{sec:retrieval-results} \\
    \addlinespace
    \multicolumn{3}{@{}l}{\textit{Historical baselines re-scored ($n = 98$; LoRA at 73)}} \\
    Best system (RAG, $k = 5$), parent-level citation & 0.625 & Table~\ref{tab:baselines} \\
    The same system, exact citation & 0.443 & Table~\ref{tab:baselines} \\
    The same system on the 25 records the draft excluded & 0.360 & Table~\ref{tab:baseline-decomposition} \\
    Few-shot $k = 10$, original scorer then exact & 0.286 / 0.159 & Table~\ref{tab:baselines} \\
    \addlinespace
    \multicolumn{3}{@{}l}{\textit{Evidence-constrained generation, 140 answers}} \\
    Quote verification passed, dev / held-out & 95.2 / 87.5\% & Table~\ref{tab:prompt-ablation} \\
    Answers containing an elided quote, dev / held-out & 0.0 / 7.1\% & Table~\ref{tab:prompt-ablation} \\
    Evidence selection correct, dev / held-out & 78.6 / 83.9\% & Table~\ref{tab:prompt-ablation} \\
    Schema contract satisfied, dev / held-out & 100 / 98.2\% & Table~\ref{tab:prompt-ablation} \\
    Abstentions & 0 of 140 & \S\ref{sec:no-abstention} \\
    Enriched context: evidence selection, input tokens & 81.0\% / $+$79\% & Table~\ref{tab:context-variants} \\
    \addlinespace
    \multicolumn{3}{@{}l}{\textit{Escalation}} \\
    Development error base rate & 21.4\% (9 of 42) & \S\ref{sec:escalation} \\
    Best candidate rule, precision [95\% CI] & 0.42 [0.23, 0.64] & Table~\ref{tab:escalation} \\
    Questions required to demonstrate it & $\approx 178$ & Table~\ref{tab:power} \\
    Answers failing in both modes & 0 & Figure~\ref{fig:failure-modes} \\
    \bottomrule
  \end{tabularx}
\end{table}

Three results are stated in advance because they shape what follows; none of
them is a performance number.

\begin{enumerate}
  \item Four aviation professionals shown the same benchmark records agreed on
    the overall verdict for fewer than half of them
    (Section~\ref{sec:review-results}). A single gold label was never as certain
    as a single number implied.
  \item The system's two failure modes --- selecting the wrong provision, and
    misreporting what the selected provision says --- did not co-occur even once
    on the development set (Section~\ref{sec:escalation}). They are independent,
    and detecting one tells the system nothing about the other.
  \item Given a schema field, an explicit instruction, and no cost for using it,
    the generator declined to abstain on every one of 140 opportunities across
    both partitions, including
    on answers that were demonstrably ungrounded
    (Section~\ref{sec:no-abstention}). Abstention cannot be delegated to the
    generator's own judgement.
\end{enumerate}

\section{The Corpus and the Frozen Benchmark}
\label{sec:frozen-benchmark}

\subsection{Structural corpus}

The parent--child corpus contains 109 legal parents and 1{,}535 retrieval
children, drawn from Regulation (EU) No 376/2014 together with its implementing
and delegated instruments and the associated Guidance
Material~\cite{eu2014reg376,easa2022easyaccess376}. The build is deterministic:
rebuilding from the same source export reproduces the corpus byte for byte, and
the structural validation records zero issues. Short units are preserved rather
than merged, because a two-line definition is a meaningful legal object and blind
minimum-length merging would reintroduce the boundary confounding that the corpus
exists to remove.

\subsection{The multi-reviewer validity gate}
\label{sec:review-results}

The benchmark was AI-assisted in construction and AI-audited for evidence
linkage. That is not a validity claim, and no amount of internal consistency
makes it one. A bounded independent review was therefore run: 33 records --- every
audit exception plus a deterministic stratified sample of unchanged records ---
were put to practising aviation professionals through a bilingual form, under an
aggregation rule fixed before any response arrived.

Four reviewers responded. Table~\ref{tab:reviewers} describes them.

\begin{table}[htbp]
  \centering
  \caption{Independent review panel. Reviewers are pseudonymised; the mapping to identities is held separately and is not part of this work.}
  \label{tab:reviewers}
  \begin{tabular}{llll}
    \toprule
    \textbf{Reviewer} & \textbf{Records rated} & \textbf{Experience} & \textbf{Familiarity with 376/2014} \\
    \midrule
    R1 & 22 & over 10 years & knows the regulation \\
    R2 & 12 & over 10 years & works with it directly \\
    R3 & 29 & 2--5 years & works with it directly \\
    R4 & 33 & over 10 years & works with it directly \\
    \bottomrule
  \end{tabular}
\end{table}

All 33 records received at least one rating and 31 received at least two.
Agreement on the overall verdict is reported in Table~\ref{tab:agreement}.

\begin{table}[htbp]
  \centering
  \caption{Inter-rater agreement on the overall verdict, computed over the 31 records with at least two raters.}
  \label{tab:agreement}
  \begin{tabular}{lr}
    \toprule
    \textbf{Quantity} & \textbf{Value} \\
    \midrule
    Records with $\geq 2$ raters & 31 \\
    Exact agreement on the verdict & 13/31 = 41.9\% \\
    Fleiss' $\kappa$, three categories & 0.250 \\
    Fleiss' $\kappa$, accept versus needs-change & 0.332 \\
    \bottomrule
  \end{tabular}
\end{table}

Both $\kappa$ values fall in the \enquote{fair} band of the Landis and Koch
scale~\cite{landis1977measurement}, computed with the variable-rater
generalisation of Fleiss' statistic~\cite{fleiss1971measuring,gwet2014interrater} because coverage
per record is uneven. The second value collapses the three verdict categories
into a binary accept-versus-change distinction. That collapse was decided
\emph{after} seeing the data and is therefore reported as a post-hoc
transformation, not as the headline.

The 41.9\% figure could be read as a failure of the review. That is not the
reading the data supports. Reviewer severity differed substantially --- individual acceptance rates
ranged from 33\% to 69\%, and two reviewers never once voted to exclude a record.
The disagreement is concentrated on precisely the records that the internal audit
had already flagged as difficult. The reading that follows is the less comfortable one:
\textbf{if four professionals who work with this regulation cannot agree whether
an answer is supported by a provision, then a single gold label for those records
was never as certain as a single accuracy figure suggested}. This is the
empirical foundation for the abstention argument developed in
Section~\ref{sec:escalation}, and it is an argument this work can make only
because the review was run.

\subsection{What the review changed}

Fourteen records were cleared by the panel under the pre-registered rule. The
remaining 19 were adjudicated by the author, each with a recorded one-line
justification, and reported separately from the panel-cleared set. The freeze
outcome is given in Table~\ref{tab:freeze}.

\begin{table}[htbp]
  \centering
  \caption{Benchmark freeze outcome.}
  \label{tab:freeze}
  \begin{tabular}{lr}
    \toprule
    \textbf{Quantity} & \textbf{Value} \\
    \midrule
    Records entering review & 33 \\
    \quad accepted & 19 \\
    \quad corrected & 12 \\
    \quad excluded & 2 \\
    \midrule
    Benchmark before freeze & 99 \\
    Benchmark after freeze & \textbf{98} \\
    Development / held-out partition & 42 / 56 \\
    Distinct legal parents & 60 \\
    Parents appearing in both partitions & \textbf{0} \\
    \bottomrule
  \end{tabular}
\end{table}

Two records were excluded. \texttt{QA\_0003} was already marked invalid by the
audit and the panel confirmed it. \texttt{QA\_0063} --- one of the three residual
retrieval failures identified before the review --- was accepted by no reviewer at
all, splitting two-to-two between \enquote{correct} and \enquote{exclude}. The
prior hypothesis that this record was defective rather than a system failure was
therefore confirmed from outside the system.

The two other residuals moved in the opposite direction. \texttt{QA\_0080} and
\texttt{QA\_0176} were both \emph{accepted} by the panel, which means the
retrieval misses on those records are genuine system failures and cannot be
explained away as incomplete gold labels. An earlier internal analysis had argued
that \texttt{QA\_0080} had equivalent evidence elsewhere in the annex; the panel
disagreed, and the panel's verdict stands.

One methodological point deserves emphasis because it was nearly missed. The
freeze script originally recomputed the parent-grouped partition from the
retained records. Because that procedure fills groups greedily toward a target
proportion, removing a single record reordered the groups and moved \textbf{12
records across the development/held-out boundary}. Every earlier selection
decision --- the embedding model, the fusion configuration, the reranker --- had
been made on the previous development set, so records used for selection would
have landed in the held-out partition. The partition is therefore \emph{inherited}
rather than recomputed: no record changed side, and every earlier selection
remains valid. The observed development fraction is consequently 42/98 = 0.429
rather than the nominal 0.400, and it is reported as observed.

\section{Retrieval}
\label{sec:retrieval-results}

All figures in this section are computed against the frozen reviewed labels on
the inherited partition. Selection decisions were made on development data;
held-out figures are reported after the fact and were not used to choose anything.

\subsection{Dense embedding comparison}

Both embedding models index identical children and identical questions.
Table~\ref{tab:dense-dev} reports the development criterion on which the
selection was made.

\begin{table}[htbp]
  \centering
  \caption{Development recall for the dense embedding comparison, reviewed labels. Proportion of questions retrieving at least one acceptable evidence child.}
  \label{tab:dense-dev}
  \begin{tabular}{lrrr}
    \toprule
    \textbf{Dense model} & \textbf{@5} & \textbf{@10} & \textbf{nDCG@10} \\
    \midrule
    \texttt{text-embedding-3-small} & 66.7\% & \textbf{78.6\%} & 0.504 \\
    \texttt{text-embedding-3-large} & \textbf{71.4\%} & 76.2\% & \textbf{0.536} \\
    \bottomrule
  \end{tabular}
\end{table}

The larger model was selected. On reviewed development labels it leads on five of
the six criteria examined --- recall and nDCG at depth 5, evidence recall at
depths 5 and 10, and nDCG at depth 10 --- but it \emph{loses} recall at depth 10,
78.6\% against 76.2\%. Under the pre-review labels it led at both depths.

This reversal is reported rather than smoothed over, and the selection is not
re-made. Re-choosing the embedding model after seeing reviewed labels would make
the choice depend on the same data that grounds the reported results. The
operating point of the deployed configuration is depth 5 followed by reranking,
where the larger model still leads. The honest summary is that the two embedding
models are closer than the original margin suggested, and that a selection
justified on a single cutoff would have been fragile.

\subsection{Lexical, hybrid and reranked retrieval}

Table~\ref{tab:retrieval-test} reports held-out performance for all five
configurations, and Figure~\ref{fig:reviewed-retrieval} shows the same data.

\begin{table}[htbp]
  \centering
  \caption{Held-out retrieval on the frozen reviewed benchmark ($n = 56$). Proportion of questions retrieving at least one acceptable evidence child.}
  \label{tab:retrieval-test}
  \rowcolors{2}{aivancityLight}{white}
  \begin{tabular}{lrrrrr}
    \toprule
    \textbf{Method} & \textbf{@1} & \textbf{@3} & \textbf{@5} & \textbf{@10} & \textbf{@20} \\
    \midrule
    Dense \texttt{3-small} & 51.8\% & 60.7\% & 67.9\% & 80.4\% & 91.1\% \\
    Dense \texttt{3-large} & 46.4\% & 64.3\% & 76.8\% & 83.9\% & 89.3\% \\
    BM25 & 53.6\% & 71.4\% & 75.0\% & 85.7\% & 85.7\% \\
    Hybrid RRF & 55.4\% & 78.6\% & 83.9\% & 91.1\% & \textbf{96.4\%} \\
    Hybrid + cross-encoder & \textbf{71.4\%} & \textbf{89.3\%} & \textbf{89.3\%} & \textbf{94.6\%} & \textbf{96.4\%} \\
    \bottomrule
  \end{tabular}
\end{table}

\begin{figure}[htbp]
  \centering
  \includegraphics[width=0.95\textwidth]{reviewed_retrieval_comparison.png}
  \caption[Held-out retrieval on the frozen benchmark]{Held-out retrieval on the frozen, independently reviewed benchmark. Left: recall by cutoff for all five configurations. Right: reranking changes the order of the candidate pool without changing its ceiling. Source: stored experiment metrics files.}
  \label{fig:reviewed-retrieval}
\end{figure}

Three observations follow.

\textbf{The channels are complementary.} BM25 is strong at exactly the places
dense retrieval is weak --- article numbers, defined terms, structured
cross-references --- and dense retrieval recovers paraphrased questions that share
no vocabulary with the provision. Equal-weight reciprocal-rank
fusion~\cite{cormack2009rrf} improves on both at every depth. It was chosen
because an untuned fusion function has no parameters that could leak information
from the held-out partition. This is a deliberate trade of performance for
methodological safety, and it should be stated as such: convex combination of
normalised scores is known to outperform reciprocal-rank fusion when its weights
are tuned~\cite{bruch2024fusion}, and tuning them here would have required
fitting on data that also grounds the reported result.

\textbf{Reranking reorders; it does not enlarge.} The cross-encoder raises top-1
recall from 55.4\% to 71.4\% and depth-5 recall from 83.9\% to 89.3\%, while
depth-20 recall is identical at 96.4\%. This is the expected signature of a
reranking stage~\cite{nogueira2019reranking}: the candidate pool is fixed by
fusion, and the reranker only decides what reaches a small generation budget. Any
evidence absent from the top 20 is unrecoverable downstream, which is why the
escalation policy in Section~\ref{sec:escalation} treats pool exhaustion as a
distinct condition from mis-ordering.

\textbf{Evidence recall trails child recall.} Under the mean-evidence-recall
measure --- the proportion of a record's required evidence set that is retrieved,
rather than whether any of it is --- the reranked configuration reaches 81.1\% at
depth 5 and 85.6\% at depth 10 on held-out data. The gap between this and the
89.3\% single-child figure is entirely accounted for by records whose gold
evidence spans several children of one legal parent.

\subsection{Effect of the review on measured performance}

Because the partition was inherited, the pre-review and post-review runs are
directly comparable: same records, same split, one record fewer in the held-out
set. Table~\ref{tab:before-after} gives the comparison for the deployed
configuration.

\begin{table}[htbp]
  \centering
  \caption{Held-out recall for hybrid retrieval with cross-encoder reranking, before and after independent review.}
  \label{tab:before-after}
  \begin{tabular}{lrrrr}
    \toprule
    \textbf{Labels} & \textbf{@1} & \textbf{@5} & \textbf{@10} & \textbf{@20} \\
    \midrule
    AI-assisted, pre-review ($n = 57$) & 66.7\% & 87.7\% & 93.0\% & 94.7\% \\
    Frozen, reviewed ($n = 56$) & \textbf{71.4\%} & \textbf{89.3\%} & \textbf{94.6\%} & \textbf{96.4\%} \\
    \bottomrule
  \end{tabular}
\end{table}

Every system improved, by one to five points. This direction matters for the
interpretation of the review: the corrections mostly repaired labels that had been
marking correct retrievals as misses, rather than exposing failures the earlier
labels had concealed. The review could have been presented as having uncovered hidden weakness. It
did not, and the direction is reported as observed.

\subsection{Rejected alternatives}

Two simpler interventions were tested on development data and rejected.
Parent-mass boosting, which adds accumulated scores from sibling children of the
same legal parent, consistently reduced recall: long or repetitive guidance
sections accumulate score without any guarantee that the specific provision is
relevant. Unequal fusion weights likewise failed to improve the top-20 candidate
objective; a BM25-heavy setting recovered one development question at depth 5
while reducing deeper recall.

Both are retained in the record as negative ablations. They establish that the
gain from the cross-encoder comes from direct query--passage interaction rather
than from hierarchy awareness in general, and that structural knowledge of the
corpus is not automatically useful just because the corpus is structured.

\section{The Three-System Baseline, Re-scored}
\label{sec:baselines}

The earlier comparison of few-shot prompting, retrieval-augmented generation and
a LoRA-adapted 3B model~\cite{hu2022lora} is a historical baseline rather than a
contribution of this work. It is re-scored here because its original evaluation
used a citation measure that does not survive inspection.

\subsection{What the original measure counted}

The original scorer returned a full match whenever the gold citation string
occurred anywhere in the prediction, and half credit for naming the correct
article number alone; a record counted as correct at half credit or above,
combined with a lexical content threshold. Both rules are unsound. Substring
containment means a gold citation of \enquote{Article 4} matches inside a
prediction that says \enquote{Article 4(6)}, and \enquote{Section 1} matches
inside \enquote{section 1.2}. More seriously, half credit for the article number
passes the correctness threshold, so a prediction citing Article 16(1) where the
provision is Article 16(6) was recorded as correct. In citation-grounded legal
question answering that is not a near miss; it is the wrong provision, and it is
the specific failure the system exists to prevent.

The replacement resolves citations against the corpus rather than comparing
strings. Two quantities are reported. \emph{Parent-level} correctness holds when
the prediction names a citation that resolves, against the parent corpus, to the
provision the gold citation denotes; only identifiers that exist in the corpus can
be credited, so an invented reference resolves to nothing. \emph{Exact}
correctness additionally requires every atomic unit of the gold citation,
including the sub-paragraph, to be named. Where a gold citation does not state
which regulation it means, every regulation containing that provision counts as a
candidate: the benchmark's own label does not make the distinction, and a system
is not penalised for reproducing an ambiguity it inherited.

\subsection{Scope}

Of the 98 frozen records, 25 carry a question the stored predictions never saw
--- 8 reworded during the independent review and 17 reworded earlier by the AI
audit. Scoring a stored answer against a question that was never put to the model
measures nothing. An earlier draft of this chapter therefore reported the
comparison over the 73 records where a stored prediction was still a valid
answer, and warned that the exclusion was not neutral: the audit rewrote the
questions it judged vague or compound, so the surviving 73 were plausibly the
better-posed subset.

That warning is now a measurement rather than a caveat. The few-shot and
retrieval-augmented predictions were regenerated against the frozen wording over
all 98 records, with the historical system prompts reproduced verbatim, the same
exemplar-sampling procedure at seed 42, temperature 0, the same $k$ values and
the same legacy flat chunk index. Two things differ from the original run and are
recorded rather than hidden: the model is pinned to
\texttt{gpt-4o-mini-2024-07-18} instead of the floating \texttt{gpt-4o-mini}
alias, so the comparison cannot change model underneath itself; and retrieval
uses NumPy cosine similarity over L2-normalised vectors in place of
\texttt{faiss.IndexFlatIP}, which computes the same quantity with ties broken by
chunk order (decision D049). The comparison below therefore covers \textbf{all 98
records}, of which 88 have a gold citation that resolves to a provision at all.

The LoRA system is the exception and is reported at $n = 73$. Re-running it needs
the base model weights and a local GPU rather than an API call, and it was not
re-run before submission; its stored predictions still answer the pre-freeze
wording, so the scorer drops the same 25 records for it. Every table below carries
its denominator, and the asymmetry is stated wherever the systems are compared.

\subsection{Results}

\begin{table}[htbp]
  \centering
  \caption{Historical baselines re-scored on the frozen benchmark. Few-shot and
  RAG were regenerated against the frozen question wording and cover all 98
  records, with citation measures over the 88 whose gold citation resolves to a
  provision. LoRA retains its stored predictions and is therefore scored on the 73
  records whose wording did not change, 63 of them resolvable. The loose measure
  is the original scorer, reproduced verbatim and reported for continuity only.}
  \label{tab:baselines}
  \rowcolors{2}{aivancityLight}{white}
  \begin{tabular}{lrrrrr}
    \toprule
    \textbf{System} & \textbf{$n$} & \textbf{Loose (original)} & \textbf{Parent-level} & \textbf{Exact} & \textbf{ROUGE-L} \\
    \midrule
    Few-shot, $k=1$ & 98 & 0.235 & 0.318 & 0.193 & 0.246 \\
    Few-shot, $k=5$ & 98 & 0.286 & 0.330 & 0.159 & 0.239 \\
    Few-shot, $k=10$ & 98 & 0.286 & 0.318 & 0.159 & 0.246 \\
    LoRA 3B & 73 & 0.178 & 0.238 & 0.111 & 0.277 \\
    RAG, $k=1$ & 98 & 0.429 & 0.409 & 0.330 & 0.338 \\
    RAG, $k=3$ & 98 & 0.592 & 0.568 & 0.420 & 0.410 \\
    RAG, $k=5$ & 98 & \textbf{0.633} & \textbf{0.625} & \textbf{0.443} & \textbf{0.432} \\
    \bottomrule
  \end{tabular}
\end{table}

\begin{table}[htbp]
  \centering
  \caption{Where the change from the earlier draft's figures comes from, at
  parent level. \emph{Stored, 73} is the draft's comparison. \emph{Regenerated,
  73} holds the record set fixed and changes only the predictions, so it measures
  run-to-run variation. \emph{All 98} holds the provenance of the predictions
  fixed and adds the 25 records the draft excluded. Denominators are 63 and 88
  resolvable records.}
  \label{tab:baseline-decomposition}
  \small
  \begin{tabular}{lrrr}
    \toprule
    & \multicolumn{2}{c}{\textbf{Same 73 records}} & \\
    \cmidrule(lr){2-3}
    \textbf{System} & \textbf{Stored} & \textbf{Regenerated} & \textbf{All 98} \\
    \midrule
    Few-shot, $k=1$ & 0.270 & 0.238 & 0.318 \\
    Few-shot, $k=5$ & 0.254 & 0.302 & 0.330 \\
    Few-shot, $k=10$ & 0.302 & 0.270 & 0.318 \\
    RAG, $k=1$ & 0.492 & 0.492 & 0.409 \\
    RAG, $k=3$ & 0.667 & 0.683 & 0.568 \\
    RAG, $k=5$ & \textbf{0.730} & \textbf{0.730} & \textbf{0.625} \\
    \bottomrule
  \end{tabular}
\end{table}

\textbf{The ranking is unchanged.} Retrieval-augmented generation beats few-shot
prompting, which beats the LoRA-adapted model, under every measure. The conclusion
of the earlier comparison survives both the corrected evaluator and the extension
to the full benchmark, and it is stated first.

\textbf{The re-run reproduces the stored predictions where it should.}
Table~\ref{tab:baseline-decomposition} separates the two things that changed.
Holding the record set at 73 and regenerating the predictions leaves RAG
essentially where it was: $k=1$ is identical on every measure, and $k=3$ and
$k=5$ differ by at most one record. Few-shot moves further, by up to two records,
which is expected rather than troubling --- its exemplars are drawn at random and
the draw sequence changes when the record set does. The middle column is a
run-to-run control, and it is tight for the system the argument rests on.

\textbf{The excluded records were harder for RAG, and the draft's figures were
inflated by their absence.} On the 25 records the earlier draft dropped, RAG at
$k=5$ reaches 0.360 [0.202, 0.555] at parent level against 0.730 [0.610, 0.824]
on the 73 it kept. The intervals do not overlap. Adding those records costs the
best system ten points at parent level and eight at exact. The scope note in the
draft predicted exactly this, and the prediction is now confirmed rather than
hedged.

\textbf{They were not harder for few-shot, which points at retrieval as the
mechanism.} On the same 25 records, few-shot at $k=10$ moves the other way:
exact-citation 0.360 [0.202, 0.555] against 0.079 [0.034, 0.173] on the 73. The
questions the audit rewrote for vagueness are ones a retriever handles worse and
a model answering from parametric knowledge handles better, which is what one
would expect if the difference is mediated by retrieval --- few-shot does not
retrieve at all. On that subset the ordering of the two systems reverses, though
25 records cannot establish that the reversal is real, and it is reported as a
direction rather than a result.

\textbf{The loose measure still overstates most severely where the dominant
failure is a plausible-looking wrong citation.} Few-shot at $k=10$ falls from
0.286 under the original scorer to 0.159 exact, because its characteristic error
is to name a real article with the wrong sub-paragraph, which is what the
half-credit rule rewarded. Fourteen of its records counted correct under the
original scorer cite the wrong sub-paragraph of the right article: Article 18(5)
where the provision is Article 18(2), Article 16(1) where it is Article 16(6).

\textbf{The most informative quantity is the gap within the RAG column.} At
$k=5$, retrieval-augmented generation reaches 0.625 [0.521, 0.719] at parent level
but 0.443 [0.344, 0.547] exact. Those 18 points are sixteen answers that retrieved
the correct article, read fluently, and cited the wrong paragraph of it. The gap
was 22 points on the draft's 73 records and it is 18 on the full benchmark: the
correction narrows it without touching the finding. No metric based on substring
citation matching can see that gap, and no lexical content measure detects it
either. It is the empirical motivation for the verification layer.

ROUGE-L~\cite{lin2004rouge} is reported because the original evaluation used it
and continuity matters. It is a lexical proxy. It is not evidence of semantic
support, and no threshold on it is described in this work as correctness.

\section{Evidence-Constrained Generation}
\label{sec:generation}

Generation was run on the development partition ($n = 42$) and then once on the
held-out partition ($n = 56$), using \texttt{gpt-4o-mini-2024-07-18} at
temperature 0. The held-out run was made only after the prompt, the verification
rule, the retrieval stack and the benchmark were all fixed, and no component was
revised after seeing its result. Both partitions are reported together
throughout; where they disagree, they are shown disagreeing. The generator is held
fixed at the model used for the historical few-shot and RAG baselines, so that
the comparison between this pipeline and those baselines isolates the scaffolding
--- retrieval quality, the evidence constraint and the output contract --- rather
than model capability. A more capable generator would also suppress the failures
this chapter exists to study.

\subsection{A schema-conformant answer is not a contract-valid answer}
\label{sec:contract-vs-schema}

Across 140 generated answers there were zero parse failures and zero refusals.
The request uses the API's strict structured-output mode~\cite{openai2024structured}, derived from the
project's own schema, and every response is then re-validated locally against
that schema. On the development partition nothing failed. On the held-out partition, one answer did, and that case bears directly on the
argument of this chapter.

For \texttt{QA\_0017} the model attached a claim to a source identifier it was
never given:

\begin{center}
\small
\begin{tabular}{ll}
  supplied & \texttt{EU\_376\_2014\_ART\_5\_\_C0001} \\
  emitted  & \texttt{EU\_376\_2014\_ART\_5\_\_C001} \\
\end{tabular}
\end{center}

\noindent The reference is fabricated, and it differs from a real one by a single
omitted character.

The API's strict mode had already accepted this response, and correctly so.
Strict structured output guarantees the \emph{shape} of a reply: that the field
exists and contains a string. It cannot guarantee that the string is one the
application supplied, because that is not a property of the schema --- it is a
property of the request. The distinction is invisible in normal operation and
becomes visible exactly when a model produces a plausible near-miss, which is
also precisely the case a human reviewer skimming the answer would not catch.
The deterministic allowed-source check caught it immediately.

Two consequences follow. Structural conformance is otherwise not a confound:
139 of 140 answers were well formed, so the failures reported below are failures
of grounding rather than of format. And provider-side structured output does not
remove the need for application-side validation of the values it carries; the
architecture keeps local re-validation after every generation call regardless of
what the provider guarantees.

\subsection{What counts as a verbatim quote}

The contract requires every claim to carry a source identifier and a verbatim
supporting quote, and forbids the model from writing display citations at all ---
those are resolved by the application from the identifier. Verification of the
quote is deterministic, but the definition of \enquote{verbatim} is a design
decision that moves the reported grounding rate by twenty-one points, so it is
stated explicitly.

Of 132 quote spans in the first run, 90 were byte-exact. \textbf{Zero} were
repaired by normalising unicode or typographic quotation marks --- so typography
is not what breaks these quotes. Thirty-one contained exactly the same words in
the same order and differed only in punctuation, almost all being the model ending
a quote early and closing it with a full stop where the source continued with a
comma. Eleven failed genuinely.

A quote is therefore accepted when its word sequence appears \emph{contiguously}
in the cited source after punctuation and typographic normalisation. Punctuation
is ignored because the corpus is a PDF extraction that carries its own artefacts
--- \texttt{safetyrelated} appears unhyphenated in the source text --- so requiring
byte-identical punctuation partly measures the extractor. Contiguity is \emph{not} relaxed, and it is the condition that does the
discriminating: it is what rejects a quote assembled from fragments. The byte-level result is reported alongside
throughout as a second figure.

\subsection{One prompt rule, measured}

The first prompt asked for a \enquote{verbatim supporting quote} without saying
what that excluded. It produced elided quotes --- quotes with material omitted
behind an ellipsis --- in 7 of 42 answers. In every case the omission removed
part of a legal test. On \texttt{QA\_0012}, the model quoted a provision on
unmanned aircraft as
\enquote{\ldots not required\ldots{} unless the occurrence\ldots{} resulted in a
fatal or serious injury to a person}, deleting both the cross-reference
identifying which aircraft the rule covers and an entire alternative limb of the
exception. The resulting answer is fluent, cites the correct article, and states
the exception more narrowly than the law does.

A second prompt version changes exactly one rule: quote one unbroken run of
consecutive words, never use an ellipsis, never join two passages, and if the
required span is long then quote less or provide several separate supports.
Nothing else differs. Table~\ref{tab:prompt-ablation} reports the effect.

\begin{table}[htbp]
  \centering
  \caption{Effect of a single prompt rule forbidding elided quotes, and how it transfers to unseen questions.}
  \label{tab:prompt-ablation}
  \small
  \rowcolors{2}{aivancityLight}{white}
  \begin{tabular}{lrrr}
    \toprule
    & \multicolumn{2}{c}{\textbf{Development} ($n=42$)} & \textbf{Held-out} ($n=56$) \\
    \cmidrule(lr){2-3}\cmidrule(lr){4-4}
    \textbf{Quantity} & \textbf{v1} & \textbf{v2} & \textbf{v2} \\
    \midrule
    Quote verification passed & 76.2\% & \textbf{95.2\%} & 87.5\% \\
    \quad under the byte-level rule & 54.8\% & 73.8\% & 64.3\% \\
    Answers containing an elided quote & 16.7\% & \textbf{0.0\%} & \textbf{7.1\%} \\
    Evidence selection correct & 81.0\% & 78.6\% & \textbf{83.9\%} \\
    Schema contract satisfied & 100\% & 100\% & 98.2\% \\
    Model abstained & 0 & 0 & 0 \\
    \bottomrule
  \end{tabular}
\end{table}

Nine development records were repaired for the cost of one sentence of prompt.
The claim count fell from 124 to 114 while the number of quote spans held steady
at 132 and 133, which suggests the constraint suppresses claims the model could
not ground rather than degrading the answers.

\textbf{The rule did not transfer as cleanly as development suggested.} On
development it eliminated elision entirely --- 7 answers to 0. On held-out
questions it did not: 4 of 56 answers still contain an elided quote. The development zero was optimistic, and the statement the data supports is
that the rule substantially reduced the failure without removing it. Grounding fell from 95.2\%
to 87.5\%; the Wilson intervals overlap ([0.842, 0.987] against [0.764, 0.938]),
so that particular drop is not distinguishable from sampling variation at this
size, but the reappearance of elisions is a qualitative recurrence rather than a
question of interval width.

The simplest intervention that addresses the failure was tried first and
measured rather than assumed, and measuring it on unseen data is what showed
that a fix validated on development can look complete and not be. That is an argument for keeping the
verifier in place permanently rather than trusting the instruction that was meant
to make it unnecessary.

\subsection{The failure the fix displaced}

The second prompt removed every \emph{marked} elision. It introduced one
\emph{unmarked} one.

On \texttt{QA\_0294}, instructed not to write an ellipsis, the model did not
quote a shorter continuous span. It joined two fragments with no marker at all,
producing \enquote{who have the role to verify compliance with applicable design
data\ldots} from a source reading \enquote{\ldots who are directly involved in the
production of aeronautical items, have the role to verify compliance with
applicable design data\ldots}. The omitted clause is the one that limits the
provision to persons directly involved in production.

The consequence is methodological. A syntactic check for ellipsis characters
would have passed this answer; only verifying that the word sequence actually
appears contiguously in the source rejects it. \textbf{A prompt constraint
displaced the failure mode into a less visible form rather than removing it.}
This is the argument for verification over instruction, and it is the argument
this project makes generally: a rule stated to a generator is a request, whereas
a rule checked against the evidence is a guarantee.

The one other surviving failure is plainer. On \texttt{QA\_0047} the model
asserted that \enquote{Member States and EASA are not allowed to record personal
details in their databases}, attributing it to a source that says only that
organisations must keep occurrence details confidential. The same structural check
rejects it for the same reason.

\subsection{The enriched context variant}
\label{sec:enriched-context}

Chapter~\ref{chap:method} specifies two context variants and everything reported
so far uses the compact one: the five reranked children under a 1{,}200-word
budget. The enriched variant keeps the same five anchors and adds their sibling
units at radius one. It is evaluated here on the development partition only. A
second post-freeze run on the held-out partition would retire the single-run discipline established at the start of this section, and an ablation is not worth that (decision D050).

\begin{table}[htbp]
  \centering
  \caption{Compact against enriched context on the development partition ($n = 42$). Prompt, schema, model, temperature and question set are identical; only the context construction differs.}
  \label{tab:context-variants}
  \small
  \rowcolors{2}{aivancityLight}{white}
  \begin{tabular}{lrr}
    \toprule
    \textbf{Quantity} & \textbf{Compact} & \textbf{Enriched} \\
    \midrule
    Quote verification passed & 40 (95.2\%) & 40 (95.2\%) \\
    \quad under the byte-level rule & 31 (73.8\%) & 32 (76.2\%) \\
    Answers containing an elided quote & \textbf{0} & 1 \\
    Evidence selection correct & 33 (78.6\%) & 34 (81.0\%) \\
    Schema contract satisfied & 42 (100\%) & 42 (100\%) \\
    Model abstained & 0 & 0 \\
    \midrule
    Claims made & 114 & 126 \\
    Quote spans & 133 & 140 \\
    Distinct sources cited per answer (mean) & 2.60 & 2.98 \\
    \midrule
    Prompt tokens & 70{,}089 & \textbf{125{,}205} \\
    Completion tokens & 13{,}747 & 14{,}673 \\
    \bottomrule
  \end{tabular}
\end{table}

The enriched context costs 79\% more input tokens and returns one record on
evidence selection --- 34 of 42 against 33, or 0.810 [0.667, 0.900] against 0.786
[0.641, 0.883] --- and one on byte-exact quoting. It also reintroduces an elided
quote where the compact run had none. At this sample size those movements are
each a single answer, and the Wilson intervals overlap almost entirely; the only
difference measured with any confidence is the cost.

One pattern is worth naming because it is not a matter of interval width. The
enriched run makes more claims (126 against 114), cites more distinct sources per
answer (2.98 against 2.60) and produces more quote spans (140 against 133). More
evidence in the window produced more assertions rather than better-grounded ones:
the grounding rate is identical at 95.2\% and the elided quote reappeared. This is
consistent with the reason the budget exists at all --- evidence is used
unreliably as the window grows~\cite{liu2024lostmiddle} --- but it is one
development run and is offered as an observation, not a finding.

The variant is therefore specified, executed and reported, and it is not adopted.
The compact context remains the configuration on which the held-out result stands.

\subsection{The generator never abstains}
\label{sec:no-abstention}

Across all three runs --- 140 answers, on both partitions --- the model
abstained \textbf{zero} times.

This is not an oversight in the design. The output schema provides four
abstention reasons; the prompt instructs the model to use them when the evidence
is insufficient, ambiguous, conflicting or outside the corpus; strict structured
output guarantees that returning one is well-formed and costs nothing. It did not do so on the record where it fabricated a prohibition, nor on the
record where it spliced away a limiting clause, nor on the record where it
invented a source identifier one character away from a real one, nor on any of
the eighteen records across both partitions where it attached its claims to the
wrong provision.

The literature on selective prediction assumes a model can be induced to express
when it should not answer, whether through calibration of its own
confidence~\cite{kadavath2022know,jiang2021calibration} or through selective
answering mechanisms~\cite{kamath2020selective,geifman2017selective}. On this
task, with this model, the explicit and freely available channel for doing so was
never used. \textbf{Abstention therefore cannot be delegated to the generator's
judgement; it must be driven by external verification.} That converts the
abstention component of the architecture from a design preference into a measured
necessity.

\section{Escalation Calibration}
\label{sec:escalation}

The escalation policy asks which \emph{runtime-observable} condition justifies
declining to answer. A signal qualifies only if it can be computed when the
question arrives, without gold labels: retrieval scores, the verifier's output
and the shape of the generated answer all qualify; whether the answer selected
the correct provision does not, since that is the outcome being predicted.

\subsection{The sample cannot support a threshold}

On the development partition, 9 of 42 answers selected the wrong provision --- a
base error rate of 21.4\%. Table~\ref{tab:escalation} reports the candidate rules
with Wilson intervals~\cite{wilson1927probable} on precision.

\begin{table}[htbp]
  \centering
  \caption{Candidate escalation rules on the development partition ($n = 42$, 9 error events). Precision intervals are Wilson 95\%.}
  \label{tab:escalation}
  \rowcolors{2}{aivancityLight}{white}
  \begin{tabular}{lrrr}
    \toprule
    \textbf{Escalate when} & \textbf{Escalated} & \textbf{Precision} & \textbf{Recall} \\
    \midrule
    Reranker moved top-1 by $\geq 2$ ranks & 19 & 0.42 [0.23, 0.64] & \textbf{0.89} \\
    Only one claim was made & 5 & 0.60 [0.23, 0.88] & 0.33 \\
    Reranker moved top-1 by $\geq 5$ ranks & 7 & 0.43 [0.16, 0.75] & 0.33 \\
    Rerank margin $\leq 0.5$ & 14 & 0.36 [0.16, 0.61] & 0.56 \\
    Rerank margin $\leq 1.0$ & 18 & 0.33 [0.16, 0.56] & 0.67 \\
    Model cited $\leq 2$ distinct sources & 21 & 0.29 [0.14, 0.50] & 0.67 \\
    Quote verification failed & 2 & \textbf{0.00} [0.00, 0.66] & \textbf{0.00} \\
    Model abstained & 0 & --- & 0.00 \\
    \bottomrule
  \end{tabular}
\end{table}

Two rules have a lower confidence bound fractionally above the 0.214 base rate,
both at 0.23. That is a two-point margin estimated from nine error events, and one
of the two rests on five escalations in total. Neither is a basis for a threshold.

\textbf{No escalation threshold is fitted in this work.} A threshold chosen from nine events would not be distinguishable from noise at
this sample size. The policy is specified; its
retrieval-uncertainty thresholds are left explicitly open.

\subsection{How much data would be required}

Table~\ref{tab:power} gives the two-proportion sample size, at
$\alpha = 0.05$ and 80\% power, needed to distinguish a given escalation
precision from the observed base rate.

\begin{table}[htbp]
  \centering
  \caption{Sample size required to demonstrate a given escalation precision against the observed 21.4\% base error rate.}
  \label{tab:power}
  \begin{tabular}{lrr}
    \toprule
    \textbf{Target precision} & \textbf{Escalated decisions} & \textbf{Questions} \\
    \midrule
    0.35 & 172 & $\approx$ 383 \\
    0.42 (best observed) & \textbf{80} & \textbf{$\approx$ 178} \\
    0.50 & 43 & $\approx$ 96 \\
    0.60 & 25 & $\approx$ 56 \\
    \bottomrule
  \end{tabular}
\end{table}

Question counts assume the 45\% escalation rate observed for the most promising
rule. Against a frozen benchmark of 98 records, demonstrating the effect actually
observed would require roughly 1.8 times the current benchmark. This is a quantified limitation with a stated remedy, reported in place of a
threshold fitted to nine events.

\subsection{The two failure modes are independent}

The result that does hold at this sample size is structural rather than
numerical, and it is shown in Figure~\ref{fig:failure-modes}.

\begin{figure}[htbp]
  \centering
  %% Figure: the two failure modes on the development partition, one cell per
%% answered question (n = 42). Cells are grouped by outcome; position within
%% the grid carries no meaning.
%% Requires tikz with libraries arrows.meta, calc.
\begin{tikzpicture}[
  font=\footnotesize,
  cell/.style={draw=white, line width=1pt, minimum size=6.2mm, inner sep=0pt},
  clean/.style={cell, fill=aivancityLight},
  wrong/.style={cell, fill=aivancityNavy},
  quote/.style={cell, fill=aivancityGold},
  key/.style={draw=white, line width=0.8pt, minimum size=3.6mm, inner sep=0pt},
]

%% 42 answered questions: 9 wrong provision, 2 ungrounded quote, 31 neither.
\foreach \i in {0,...,41} {
  \pgfmathtruncatemacro{\col}{mod(\i,14)}
  \pgfmathtruncatemacro{\row}{div(\i,14)}
  \pgfmathtruncatemacro{\kind}{ifthenelse(\i<9, 0, ifthenelse(\i<11, 1, 2))}
  \ifnum\kind=0
    \node[wrong] at (\col*6.6mm, -\row*6.6mm) {};
  \else
    \ifnum\kind=1
      \node[quote] at (\col*6.6mm, -\row*6.6mm) {};
    \else
      \node[clean] at (\col*6.6mm, -\row*6.6mm) {};
    \fi
  \fi
}

%% brace over the two failing groups
\draw[draw=aivancityGray, line width=0.5pt]
  (-3.6mm, 4.0mm) -- (69.9mm, 4.0mm);
\node[anchor=south, font=\footnotesize, text=aivancityGray] at (33.0mm, 4.3mm)
  {42 answered questions, development partition};

%% ------------------------------------------------------------------- key ---
\node[key, fill=aivancityNavy, anchor=west] (kw) at (-3.6mm, -20.0mm) {};
\node[anchor=west, font=\footnotesize] at ($(kw.east)+(1.6mm,0)$)
  {wrong provision selected --- 9};

\node[key, fill=aivancityGold, anchor=west] (kq) at (58.0mm, -20.0mm) {};
\node[anchor=west, font=\footnotesize] at ($(kq.east)+(1.6mm,0)$)
  {quote not grounded --- 2};

\node[key, fill=aivancityLight, anchor=west] (kc) at (-3.6mm, -26.0mm) {};
\node[anchor=west, font=\footnotesize] at ($(kc.east)+(1.6mm,0)$)
  {neither failure --- 31};

\node[anchor=west, font=\footnotesize\bfseries, text=draftRed] at (58.0mm, -26.0mm)
  {both failures --- 0};

\end{tikzpicture}

  \caption[The two failure modes do not co-occur]{The two failure modes on the
  development partition, one cell per answered question ($n = 42$). Nine answers
  selected the wrong provision and two carried a quote that does not appear in
  the source it cites; eleven answers failed in one way or the other, and none
  failed in both. Cells are grouped by outcome, so position in the grid carries
  no meaning.}
  \label{fig:failure-modes}
\end{figure}

The two sets do not intersect. Quote verification has precision 0.00 and recall
0.00 for predicting evidence-selection errors, and the two answers it rejected had
entirely correct evidence.

The finding is architectural rather than a shortfall in performance. The system fails in two
independent ways. \emph{Evidence selection} --- attaching a claim to the wrong
provision --- is not observable at runtime and must be predicted from retrieval
signals, which the preceding subsection shows cannot yet be done reliably.
\emph{Quote fidelity} --- misreporting what the selected provision says --- is
deterministically checkable at runtime with no gold labels at all. A verifier
checking only quotes would have passed all nine evidence errors; a check on
evidence selection alone would have passed both quote failures. Both layers are
required, and the requirement is measured rather than asserted.

\subsection{The policy as specified}

The escalation policy that these results support is stated as follows.

\begin{enumerate}
  \item Retrieve a broad candidate pool through lexical and dense fusion.
  \item Rerank when evidence may be present but poorly ordered.
  \item Add adjacent children when a selected provision lacks local continuity.
  \item \textbf{Block the answer when quote verification fails.} This requires no
    calibration: a quote either appears contiguously in the cited source or it
    does not, so the check has no false positives by construction. It fires on 2
    of 42 development answers.
  \item Request clarification when several legitimate provisions compete.
  \item \textbf{Escalate on retrieval uncertainty --- threshold uncalibrated.}
    Reranker-versus-fusion disagreement is the most promising observed signal and
    is named as the candidate, with its threshold left open pending the sample
    size in Table~\ref{tab:power}.
  \item Refer to a human when evidence remains absent, ambiguous or conflicting.
\end{enumerate}

Model self-abstention is excluded as a signal entirely, for the reason given in
Section~\ref{sec:no-abstention}.

\begin{figure}[htbp]
  \centering
  \resizebox{\textwidth}{!}{%% Figure: the escalation policy as a ladder of levels, with the trigger for
%% each transition and the evidential status of that trigger.
%% Requires: tikz with libraries arrows.Meta, positioning, fit, calc.
\begin{tikzpicture}[
  font=\footnotesize,
  every node/.append style={execute at begin node={%
    \hyphenpenalty=10000\exhyphenpenalty=10000\relax}},
  level box/.style={
    draw=aivancityNavy, line width=0.6pt, rounded corners=2pt,
    fill=white, align=left, inner sep=4pt, text width=6.1cm, minimum height=0.85cm
  },
  superseded/.style={
    draw=aivancityGray, line width=0.6pt, dashed, rounded corners=2pt,
    fill=white, align=left, inner sep=4pt, text width=6.1cm,
    text=aivancityGray, minimum height=0.85cm
  },
  trigger/.style={align=left, text width=5.3cm, font=\footnotesize, anchor=west},
  status/.style={align=right, text width=2.5cm, font=\footnotesize\bfseries, anchor=east},
  down/.style={-{Latex[length=4.5pt,width=4.5pt]}, draw=aivancityNavy, line width=0.8pt},
  downopen/.style={-{Latex[length=4.5pt,width=4.5pt]}, draw=draftRed, line width=0.8pt, dashed},
  downfirm/.style={-{Latex[length=4.5pt,width=4.5pt]}, draw=aivancityGold, line width=1.1pt},
]

\def\dy{1.90}

%% ------------------------------------------------------------- the ladder ---
\node[superseded]  (l0) at (0, 0*\dy) {\textbf{0}\quad Dense retrieval only};
\node[level box]   (l1) at (0,-1*\dy) {\textbf{1}\quad Hierarchy-aware retrieval};
\node[level box]   (l2) at (0,-2*\dy) {\textbf{2}\quad Hybrid candidates, rank fusion};
\node[level box]   (l3) at (0,-3*\dy) {\textbf{3}\quad Cross-encoder reranking};
\node[level box]   (l4) at (0,-4*\dy) {\textbf{4}\quad Evidence-constrained generation};
\node[level box]   (l5) at (0,-5*\dy) {\textbf{5}\quad Clarification or targeted revision};
\node[level box]   (l6) at (0,-6*\dy) {\textbf{6}\quad Abstention and human referral};

\draw[down]     (l0) -- (l1);
\draw[down]     (l1) -- (l2);
\draw[downopen] (l2) -- (l3);
\draw[down]     (l3) -- (l4);
\draw[downfirm] (l4) -- (l5);
\draw[down]     (l5) -- (l6);

%% ------------------------------------------------- triggers, right column ---
\node[trigger] at (3.35,-0.5*\dy)
  {superseded once the corpus carries its hierarchy};
\node[trigger] at (3.35,-1.5*\dy)
  {low fused score, or the lexical and dense channels disagree};
\node[trigger, text=draftRed] at (3.35,-2.5*\dy)
  {small margin between the top candidates\\
   \emph{threshold left uncalibrated}};
\node[trigger] at (3.35,-3.5*\dy)
  {applied to every question that is answered at all};
\node[trigger] at (3.35,-4.5*\dy)
  {quote verification fails, or several provisions compete legitimately};
\node[trigger] at (3.35,-5.5*\dy)
  {sources conflict, no supporting evidence, ambiguity unresolved};

%% ------------------------------------------------- status, left column ------
\node[status, text=aivancityGray] at (-3.35,-0.5*\dy) {baseline};
\node[status, text=aivancityNavy] at (-3.35,-1.5*\dy) {runtime signal};
\node[status, text=draftRed]      at (-3.35,-2.5*\dy) {open};
\node[status, text=aivancityNavy] at (-3.35,-3.5*\dy) {unconditional};
\node[status, text=aivancityGold] at (-3.35,-4.5*\dy) {deterministic};
\node[status, text=aivancityNavy] at (-3.35,-5.5*\dy) {runtime signal};

%% ------------------------------------------------------- excluded signal ---
\node[draw=aivancityGray, dashed, line width=0.6pt, rounded corners=2pt,
      fill=white, align=left, inner sep=5pt, text width=13.0cm,
      font=\footnotesize, text=aivancityGray] (excl) at (0,-6.95*\dy)
  {\textbf{Excluded as a trigger:} the generator's own abstention flag, which was never raised on any
   of the 140 opportunities across both partitions --- including on answers whose quotes did not
   appear in the sources they cited.};

%% ------------------------------------------------------------------ key ---
\node[align=center, font=\footnotesize, anchor=north, text width=13.4cm] at (0,-7.55*\dy)
  {\textcolor{aivancityGold}{\rule[0.5ex]{9pt}{1.1pt}}~deterministic, no calibration \quad
   \textcolor{aivancityNavy}{\rule[0.5ex]{9pt}{0.8pt}}~specified and applied \quad
   \textcolor{draftRed}{\rule[0.5ex]{3pt}{0.8pt}\,\rule[0.5ex]{3pt}{0.8pt}}~threshold left open};

\end{tikzpicture}
}
  \caption[The escalation policy and the status of each transition]{The
  escalation policy, with the trigger for each transition and what the results
  of this chapter establish about it. One transition needs no calibration, since
  a quote either appears contiguously in the cited source or it does not. One is
  left deliberately open: reranker-versus-fusion disagreement is the most
  promising observed signal, but nine error events cannot fix its threshold
  (Table~\ref{tab:power}). The generator's own abstention flag is excluded as a
  signal altogether.}
  \label{fig:escalation-policy}
\end{figure}

Escalation precision and recall are reported together throughout, never
separately. A rule escalating 45\% of questions to catch 89\% of errors may well
be correct for aviation regulation, where a wrong citation is expensive and a
human referral is cheap --- but that is a judgement about the relative cost of the
two outcomes, not something the data decides. Wherever an escalation rate appears,
the assumed cost ratio is stated with it.

\section{Two Properties No Internal Metric Detects}
\label{sec:invisible}

\subsection{Citation labels a reader cannot resolve}

Six records drew the same reaction from reviewers, in different words each time:
\emph{\enquote{GM 3.1 = ?}}, \emph{\enquote{qu'entends-tu par
\enquote{section 2.2}?}}. A practitioner who works with this regulation daily
could not resolve the canonical citation strings.

That observation was then measured rather than left anecdotal. Every gold citation
in the frozen benchmark was resolved against the corpus, with the result in
Table~\ref{tab:legibility}.

\begin{table}[htbp]
  \centering
  \caption{Resolution of the benchmark's own gold citation labels against the corpus ($n = 98$).}
  \label{tab:legibility}
  \begin{tabular}{lr}
    \toprule
    \textbf{The gold citation\ldots} & \textbf{Records} \\
    \midrule
    identifies the parent holding its own evidence & 48 \\
    names an Article while the evidence sits in a Guidance Material section & 28 \\
    omits which of the four regulations it means & 12 \\
    cannot be parsed into any provision at all & 10 \\
    \bottomrule
  \end{tabular}
\end{table}

\textbf{Twenty-two of 98 records carry a citation label that cannot be resolved to
a single provision from the string alone} --- roughly four times what the six
free-text comments suggested. A bare \enquote{Article 2} exists in all four
regulations in the corpus; \enquote{GM to Reg.\ (EU) No 376/2014 [04]} identifies
nothing at all.

Every one of these citations resolves correctly \emph{inside} the pipeline,
because the pipeline carries the parent identifier alongside the display string.
The defect is invisible to every internal metric by construction, and visible
immediately to an outside reader holding only what the user is actually shown.
This is the clearest available demonstration that internal consistency is not
correctness, and it was found by asking practitioners rather than by measuring.

The 28 guidance-material records are a separate and legitimate phenomenon: the
citation points to the binding provision while the evidence sits in the guidance
discussing it. It is not an error, but it means that \enquote{cited correctly} and
\enquote{retrieved the evidence} are different questions on nearly a third of the
benchmark, and results must not silently conflate them.

\subsection{The benchmark cannot demonstrate an agent's benefit}

Every one of the 98 frozen records has its gold evidence inside a single legal
parent. Zero records require evidence from two or more.

This is decisive for the question of agentic, multi-step retrieval. Whatever an
agent's merits in general, a benchmark in which no question requires combining
provisions from different parts of the instrument cannot exhibit them: the
measurement has no room to show a difference. Building an agent and reporting a
null result on this benchmark would be uninformative, and reporting a positive one
would be evidence of something other than the agent's retrieval behaviour.

A cross-reference evaluation subset is therefore a prerequisite for that line of
work, not an optional refinement --- a conclusion reached from a property of the
data rather than from a position on agent architectures.

\section{Threats to Validity}
\label{sec:threats}

\textbf{Benchmark provenance.} The benchmark is synthetic in origin and may
reflect the style of the model that generated it~\cite{sainz2023contamination}.
The independent review mitigates but does not eliminate this: 33 of 98 records
were reviewed, and the remaining 65 retain AI-assisted evidence linkage.

\textbf{Agreement is fair, not strong.} Fleiss' $\kappa$ of 0.250 means the
reviewed labels themselves carry disagreement. This is reported as a finding
rather than a nuisance, but it bounds how precisely any accuracy figure computed
against those labels can be interpreted.

\textbf{Small samples.} The escalation analysis rests on 9 error events on 42
development records, and every precision figure in Section~\ref{sec:escalation} is
indicative rather than estimated. Generation is reported on both partitions, but
56 held-out records still give wide intervals: the 87.5\% grounding rate carries a
95\% interval of [0.764, 0.938].

\textbf{Selection after observation.} The second prompt version was chosen after
seeing the first version's failures on development data. That is the legitimate
use of a development partition, and the held-out run was that prompt's first and
only exposure to unseen questions. It is reported exactly as observed, including
where it is worse than development; no prompt, threshold or model was revised
afterwards.

\textbf{Baseline scope.} The few-shot and RAG comparison now covers all 98
records; LoRA remains at 73, since its stored predictions answer the pre-freeze
wording and it was not re-run. Comparisons involving LoRA are therefore made on a
subset that is measurably easier for retrieval-augmented systems, which flatters
RAG against it rather than the reverse. The LoRA system was measured at one variant, on a
3B model, against hosted models of larger and undisclosed size; it is not a
controlled comparison of adaptation methods and is not presented as one.

\textbf{Lexical proxies.} ROUGE-L appears only for continuity with the earlier
evaluation. Structural verification establishes that a claim points at supplied
evidence and reproduces its wording; it does not establish that the evidence
supports the claim. Semantic support remains unmeasured in this
work~\cite{rashkin2023attribution,liu2023verifiability}.

\textbf{Fusion configuration.} Equal-weight reciprocal-rank fusion is known to be
suboptimal relative to tuned convex combination~\cite{bruch2024fusion}. It was
chosen because an untuned function cannot leak from the held-out partition. The
reported retrieval figures are therefore a conservative estimate of what this
architecture can achieve.
